\section{群準同型}
\begin{prop}
$\phi:G\to H$を群準同型とする。
\ben
\item \tf{$\ker{\phi}$は$G$の正規部分群である。}
\item $G/\ker{\phi} \cong \im{\phi}$
\een
\end{prop}


\begin{egprob} (2013基礎数学試験No.2 一般化)
$S_n$からアーベル群$G$の全射群準同型$\phi$が存在するならば, $G=\{ 1 \}, \mb{Z}/2\mb{Z}$であることを示せ。
\end{egprob}
\begin{egans}

$|G| \geq 3$を仮定する。 $G$がアーベル群であるという仮定によって$\phi(hgh^{-1}) = \phi(g)$が成り立つことに注意する。よって, $\phi$は共役類上では一定の値を取る。すべての互換は共役であり, ある互換が$\phi$で単位元に移ると, $\phi$は自明な準同型になるので不適。よって$\phi$はすべての互換を, $G$のある位数2の元$g_2$に移す。$S_n$は互換で生成されるので, $\phi$の像は$\{ 1 , g_2\}$に含まれる。これは2元集合だから, 全射になり得ない。また, $|G| = 2$ならば符号を取る準同型$\mr{sgn}$は$\mb{Z}/2\mb{Z}$への全射準同型である。\qed
\end{egans}

\subsection{準同型定理}

\begin{egprob} (H26-基礎科目II-4)
$G = \mb{Z}/4\mb{Z}\times \mb{Z}/6\mb{Z}\times \mb{Z}/9\mb{Z}$の指数3の部分群は何個あるか.
\end{egprob}
\begin{egans}
\tf{ $\phi:G\to \mb{Z}/3\mb{Z}$という全射群準同型をすべて求めて, $\ker{\phi}$として挙がるものをすべて求める. } なぜこれでよいかというと, 任意の指数$3$の部分群$H$に対し$\psi :G/H \to \mb{Z}/3\mb{Z}$という同型が一つあり, これによって$G\to G/H \to \mb{Z}/3\mb{Z}$は$H$を核に持つ全射準同型になるからである. また, 中国剰余定理より
\[G\cong \mb{Z}/4\mb{Z}\times \mb{Z}/2\mb{Z} \times \mb{Z}/3\mb{Z} \times \mb{Z}/9\mb{Z}\]
である. $G$の元を$(x,y,z,w)$で表すとき, $\phi(1,0,0,0) = \phi(0,1,0,0) = 0$が分かる(Why?). $\phi(0,0,1,0), \phi(0,0,0,1)\in \mb{Z}/3\mb{Z}$で$\phi$は一意に決まり, 全射性よりどちらかは0ではない. このことから8個の$\phi$があると分かり, それらから$\ker{\phi}$をすべて計算すると4個異なる部分群が出てくる. よって4個である.
\end{egans}
\begin{egans} (T氏のアイデアをもとにしたより良い解答)\\
$H$を指数$3$の部分群とする. $G/H \cong \mb{Z}/3\mb{Z}$なので\ao{$3G  \subset H$である.(ここがポイント)}  よって$H$は$G/3G \cong \mb{Z}/3\mb{Z} \times \mb{Z}/3\mb{Z}$の中の指数$3$の部分群$H/3G$に対応する. それは4個ある. 実際, 指数3ということは位数3なので, 位数3の巡回元を探せばよく $(1,0)$, $(0,1)$, $(1,1)$, $(1,-1)$で生成されるものがそれらである.
\end{egans}


\begin{prob} (Liの自作問題, 難)
$G = \prod_{n=1}^{2021} \mb{Z}/n!\mb{Z}$とする. $G$の指数$2017$の部分群はいくつあるか. (2017は素数である)
\end{prob}

\begin{egprob}
群$G$が次で与えられるとき, 群準同型$\phi:G\to \mathbb{C}_{1}$をすべて求めよ\footnote{($\spadesuit$)表現論の話。$n$次元複素ベクトル空間$V$と有限群$G$に対し, 「群$G$の$n$次元(既約)表現」というのは群準同型$G\to \to GL(V)$のことであり, とくに$n=1$のときは$GL(V)\cong \mb{C}_{1}$だから, この問題は「1次元既約表現の数を決定せよ」という問題に他ならない。}。ただし, $\mb{C}_{1}$は絶対値1の複素数の集合に, 通常の積による演算を入れた群である。
\ben
\item $\mb{Z}/n\mb{Z}$
\item Kleinの四元群$K = \{ 1,(12)(34),(13)(24),(14)(23) \} \subset \mathfrak{S}_{4}$
\item $\mb{Z}/2\mb{Z}\times \mb{Z}/2\mb{Z}$
\item $\mb{Z}/2\mb{Z}\times \mb{Z}/4\mb{Z}$
\item $\maf{S}_n$ (\tf{Hint:}互換で生成できることを用いよ。)
\item (なんでも考えてみてください！)
\een
\end{egprob}
\begin{egans}
(1): $\phi(1)$の行先は1の$n$乗根であり, その行先で準同型は決定されるので$n$個。\\
(2)
\end{egans}



\subsection{群の同型と非同型}
簡単な群の同型の例を紹介するとともに, 群論的な議論を総合的に復習する。

\begin{eg} ($PSL_{2}(\mb{F}_{2}) \cong \maf{S}_{3}$, 青雪江Ex2.6.1)\\
$G=PSL_{2}(\mb{F}_{2})$とする。$G$は$SL_{2}(\mb{F}_2)$の元で定数倍で一致するものをすべて同一視することで得られるが, 各同値類は1つの元しか持たないので$SL_{2}(\mb{F}_{2}) \cong G$である。$SL_{2}(\mb{F}_{2})$の単位元ではない元は,
\beqn
\bep
1 & 1 \\
0 & 1
\enp
,
\bep
1 & 0\\
1 & 1
\enp
,
\bep
0 & 1\\
1 & 0
\enp
,
\bep
0 & 1\\
1 & 1
\enp
,
\bep
1 & 0\\
1 & 1
\enp
\eeqn
である。このうち最初の3つだけが位数が$2$である。$X$をこの3つの元の集合とする。共役作用によって$G$は$X$に作用し, したがって準同型$\phi:G\to \maf{S}_{3}$を得る。すなわち,
\[\phi(A) = (X\ni M\mapsto A^{-1}MA\in X)\]
とおくのである。$\phi$は同型$G\cong \maf{S}_{3}$を定める。これを見るには$\phi$が単射であることを見ればよいので, すべての$M\in X$で$MA = AM$であるならば$A = 1_{G}$であることを示せばよい。これは簡単に示せる。
\end{eg}

位数が同じだが同型でないことが分かる例を挙げる. しばしば, 同型でないことを論じるには\tf{元の位数}が活躍する.

\begin{egprob}
次の群は互いに同型ではないことを示せ。ただし, $\mb{Z}_{d} = \mb{Z}/d\mb{Z}$
\ben
\item $\mb{Z}_2\times \mb{Z}_2$と$\mb{Z}_{4}$
\item $S_4$と$SL_2(\mb{F}_{3})$
\een
\end{egprob}
